\section{Method}
\label{sec:method}

\subsection{Setting}

A model is trained recursively: each generation $g \in \{0,\dots,H-1\}$ trains on data drawn
from the previous generation's outputs, optionally supplemented with human-origin examples.
The quantity under allocation is a \emph{fixed lifetime budget} $B$ of human-origin
optimizer tokens. A policy decides, at each generation, how much of the remaining budget to
spend and which candidates to spend it on. Every policy in a budget-matched comparison
consumes exactly the same lifetime human-origin token count and the same total
optimizer-token count.

\subsection{Policy-visible state}

At generation $g$ a policy observes only the remaining human-token budget, the number of
generations remaining, and a monitored per-mode under-coverage score computed from a lagged
reference. It does not observe held-out test data, another policy's outputs, or any primary
treatment outcome. A partially missing mode receives score zero; if every mode is missing,
selection falls back to a seeded random candidate order.

\subsection{Treatment families}

Four budget-matched families are compared. Let $r$ denote remaining human tokens, $F$ the
number of future generations, $s_{\text{base}}$ the base per-generation spend, and
$s_{\max}$ the per-generation cap.

\paragraph{Fresh random rescue.} Spend $\min(s_{\text{base}}, r)$ each generation; shuffle
candidates under the supplied seed and select without exceeding budget.

\paragraph{Schedule-only.} Spend $s_g$ from a frozen schedule that withholds tokens through
the early generations and spends at the cap thereafter; select in seeded random order.
Timing varies; mode targeting does not.

\paragraph{Selection-only.} Spend $\min(s_{\text{base}}, r)$ each generation; score
candidates by lagged mode under-coverage, sort descending, break ties by ascending example
identifier. Mode targeting varies; the spending schedule does not.

\paragraph{Joint time-and-mode.} Let $u$ be the maximum finite monitored mode score, and let
$\tau_{\text{low}} < \tau_{\text{high}}$ be the frozen urgency thresholds. The desired spend
is
\[
d(u) = \begin{cases}
0 & u < \tau_{\text{low}}\\
s_{\text{base}} & \tau_{\text{low}} \le u < \tau_{\text{high}}\\
s_{\max} & u \ge \tau_{\text{high}}
\end{cases}
\]
and the realised spend applies a feasibility clamp,
\[
b = \mathrm{clamp}\bigl(d(u),\; \max(0,\, r - s_{\max}F),\; \min(s_{\max},\, r)\bigr).
\]
The lower bound reserves enough budget that every remaining generation can still be served
at the cap; the upper bound enforces the per-generation cap and the remaining budget.
Candidates are then sorted by descending score, ties broken by ascending example identifier.

\subsection{Why the joint policy is not two tuned baselines composed}

This is the load-bearing design question, and it is settled by an eligibility constraint
rather than by an empirical margin.

Schedule-only may alter \emph{when} tokens are spent but may not choose \emph{which} modes
receive them. Selection-only may choose modes but spends on a fixed schedule. Composing a
tuned schedule-only policy with a tuned selection-only policy yields a fixed spending
profile paired with mode-aware selection: the schedule is still chosen in advance and cannot
respond to what monitoring reports at run time.

The joint policy is the only eligible family whose spend at generation $g$ is a function of
the monitored state at generation $g$. The urgency thresholds make spending reactive --- a
generation whose worst monitored mode is well covered spends nothing, and one whose worst
mode is badly covered spends the cap. Under a fixed lifetime budget this is a reallocation,
not an increase: tokens withheld early remain available later, and the feasibility clamp
guarantees the reserve never falls below what future generations require.

Whether that reactivity helps is an empirical question this paper does not answer.

\subsection{Frozen parameters}

\begin{table}[h]
\centering
\caption{Frozen policy parameters, generated from \texttt{configs/policy/joint.json}.}
% AUTO-GENERATED by scripts/generate_method_tables.py. DO NOT EDIT VALUES BY HAND.
\begin{tabular}{lr}
Frozen parameter & Value \\
\hline
Horizon & 10 \\
Lifetime human tokens & 100 \\
Base generation spend & 10 \\
Maximum generation spend & 20 \\
Token granularity & 10 \\
Monitoring lag (generations) & 1 \\
\hline
Desired spend, $u < 0.25$ & 0 \\
Desired spend, $0.25 \le u < 0.50$ & 10 \\
Desired spend, $u \ge 0.50$ & 20 \\
\end{tabular}

\end{table}

These were fixed in the Week~2 method freeze using interfaces, fixtures, and synthetic
results. No primary treatment outcome influenced any score, schedule, seed, threshold,
exclusion rule, or interpretation.

\subsection{Separation of the treatment families}

An earlier internal finding recorded the joint and selection-only families as
observationally identical, and random as identical to schedule-only. That finding was
subsequently traced to a reverted implementation rather than to the method: with the frozen
rule restored, all four families produce distinct trajectories under the fixture simulator,
and random separates from schedule-only. This evidence is structural and fixture-based. It
establishes that the decomposition is well posed; it is not evidence of any performance
difference on real data.
